% p098_1; source PDF page 98; manually verified editable table.
\begin{tikzpicture}[x=1cm,y=-1cm,inner sep=0pt,outer sep=0pt]
\path[use as bounding box] (0,0) rectangle (16.000,7.458);
\draw[line width=.4pt] (0.000,0.000) rectangle (3.400,1.048);
\node[align=center,text width=3.140cm,font={\linespread{1}\fontsize{10}{13.50}\selectfont\bfseries}] at (1.700,0.524) {防卫过当的成立条件};
\draw[line width=.4pt] (3.400,0.000) rectangle (7.200,1.048);
\node[align=center,text width=3.540cm,font={\linespread{1}\fontsize{10}{13.50}\selectfont\bfseries}] at (5.300,0.524) {具体要求};
\draw[line width=.4pt] (7.200,0.000) rectangle (16.000,1.048);
\node[align=center,text width=8.540cm,font={\linespread{1}\fontsize{10}{13.50}\selectfont\bfseries}] at (11.600,0.524) {范例};
\draw[line width=.4pt] (0.000,1.048) rectangle (3.400,3.653);
\node[align=left,text width=3.140cm,font={\linespread{1}\fontsize{10}{13.50}\selectfont}] at (1.700,2.351) {客观要件};
\draw[line width=.4pt] (3.400,1.048) rectangle (7.200,3.653);
\node[align=left,text width=3.540cm,font={\linespread{1}\fontsize{10}{13.50}\selectfont}] at (5.300,2.351) {发生过当结果};
\draw[line width=.4pt] (7.200,1.048) rectangle (16.000,3.653);
\node[align=left,text width=8.540cm,font={\linespread{1}\fontsize{10}{13.50}\selectfont}] at (11.600,2.351) {甲轻伤害乙，乙开枪反击，虽然手段过当，但只造成甲轻伤。对此不能认定为防卫过当。也即，即使手段过当，但没有过当结果，不构成防卫过当};
\draw[line width=.4pt] (0.000,3.653) rectangle (3.400,7.458);
\node[align=left,text width=3.140cm,font={\linespread{1}\fontsize{10}{13.50}\selectfont}] at (1.700,5.556) {主观要件};
\draw[line width=.4pt] (3.400,3.653) rectangle (7.200,7.458);
\node[align=left,text width=3.540cm,font={\linespread{1}\fontsize{10}{13.50}\selectfont}] at (5.300,5.556) {对过当结果至少有过失};
\draw[line width=.4pt] (7.200,3.653) rectangle (16.000,7.458);
\node[align=left,text width=8.540cm,font={\linespread{1}\fontsize{10}{13.50}\selectfont}] at (11.600,5.556) {甲轻伤害乙，乙用木棍反击，竟致甲死亡。（1）此时，乙属于客观意义上的防卫过当。（2）到了主观要件，如果对甲的死亡有过失，则成立最终意义上的防卫过当，以过失致人死亡罪论处；如果没过失，属于意外事件，则不负刑事责任①};
\end{tikzpicture}
